% =============================================================================
% APPENDIX V: CORE-WHEN: Context Framework (Ψ-Dimensions)
% =============================================================================
% Module: BCM2_09_CTX (Context)
% Category: CORE
% Version: 1.0 (January 2026)
% Status: Final
% Language: English
%
% PURPOSE: This appendix answers the fundamental question "When does context matter?"
% It formalizes the 8 Ψ-dimensions that characterize behavioral context:
% Institutional, Social, Cognitive, Cultural, Economic, Temporal, Material, Physical.
%
% 10C POSITION: This is CORE 4 of 10 - the "WHEN" question.
%
% NOTE: This appendix replaces the former CTW references (Context Framework).
%       All CTW cross-references should now point to V.
%
% =============================================================================

% -----------------------------------------------------------------------------
% HEADER BLOCK
% -----------------------------------------------------------------------------

\begin{tcolorbox}[colback=gray!5!white, colframe=gray!75!black,
    title=Appendix Header: V CORE-WHEN]

\begin{tabular}{@{}ll@{}}
\textbf{Code:} & V \\
\textbf{Category:} & CORE (Fundamental Theory) \\
\textbf{Full Name:} & CORE-WHEN: Context Framework (Ψ-Dimensions) \\
\textbf{Module:} & BCM2\_09\_CTX (Context) \\
\textbf{Version:} & 1.0 (January 2026) \\
\textbf{Status:} & Final \\
\textbf{Language:} & English \\
\end{tabular}

\vspace{0.5em}
\textbf{Purpose:} Formalizes the \textbf{Ψ-Framework}, answering the
fourth fundamental question of the EBF: \textit{``When does context matter?''}
Establishes that behavior is context-dependent through 8 dimensions:
$\Psi = (\Psi_I, \Psi_S, \Psi_C, \Psi_K, \Psi_E, \Psi_T, \Psi_M, \Psi_F)$.

\textbf{10C Position:} This is CORE 4 (WHEN) --- modulates WHO (AAA), WHAT (C),
and HOW (B) through context sensitivity.

\textbf{Former Alias:} This appendix consolidates the former CTW (Context Framework) references.

\end{tcolorbox}

% =============================================================================
% CROSS-REFERENCE STRUCTURE
% =============================================================================

\begin{tcolorbox}[colback=purple!5!white, colframe=purple!75!black,
    title=Chapter Linkage]

\textbf{Primary Chapter:} Chapter 9: Context and the Ψ-Framework

\vspace{0.3em}
\textbf{The Fundamental Question:}
\begin{quote}
\textit{``When does context matter? How do situational factors modulate behavior and preferences?''}
\end{quote}

\textbf{Section-Level Mapping:}
\begin{itemize}[nosep]
\item Section 9.1: Context as Multiplier $\leftrightarrow$ Section~\ref{sec:when-fundamental}
\item Section 9.2: The 8 Ψ-Dimensions $\leftrightarrow$ Section~\ref{sec:when-dimensions}
\item Section 9.3: Context Modulation Rules $\leftrightarrow$ Section~\ref{sec:when-modulation}
\item Section 9.4: MACRO-MESO-MICRO Hierarchy $\leftrightarrow$ Section~\ref{sec:when-hierarchy}
\end{itemize}

\textbf{Secondary Chapters:}
\begin{itemize}[nosep]
\item Chapter 7: Complementarity --- Context affects $\gamma$ values
\item Chapter 11: Awareness --- Context affects awareness levels
\item Chapter 17-20: Interventions --- Context-specific effectiveness
\end{itemize}

\textbf{Relationship to Other COREs:}
\begin{itemize}[nosep]
\item Appendix AAA (CORE-WHO): Context varies by welfare level
\item Appendix C (CORE-WHAT): Context affects dimension weights $\alpha_d(\Psi)$
\item Appendix B (CORE-HOW): Context affects complementarity $\gamma(\Psi)$
\item Appendix AU (CORE-AWARE): Awareness depends on context
\end{itemize}

\end{tcolorbox}

% -----------------------------------------------------------------------------
% APPENDIX SCOPE BOX
% -----------------------------------------------------------------------------

\begin{tcolorbox}[
    colback=orange!5!white,
    colframe=orange!75!black,
    title={\textbf{Appendix Scope: Ziel / In-Scope / Out-of-Scope / Constraints}},
    fonttitle=\bfseries
]

\textbf{Ziel (Objective):}
\begin{itemize}[nosep]
    \item Define the complete context framework through 8 Ψ-dimensions
\end{itemize}

\textbf{In-Scope:}
\begin{itemize}[nosep]
    \item Eight Ψ-dimension definitions (I, S, C, K, E, T, M, F)
    \item Axioms PSI-1 through PSI-10
    \item MACRO-MESO-MICRO context hierarchy
    \item Context modulation functions $f(\Psi)$
\end{itemize}

\textbf{Out-of-Scope:}
\begin{itemize}[nosep]
    \item What utility dimensions exist $\rightarrow$ Appendix C (CORE-WHAT)
    \item How components interact $\rightarrow$ Appendix B (CORE-HOW)
    \item Parameter estimation for $\Psi$ $\rightarrow$ Appendix BBB (CORE-WHERE)
    \item Context data sources $\rightarrow$ BCM2 Context Database
\end{itemize}

\textbf{Constraints:}
\begin{itemize}[nosep]
    \item Context must be observable or inferable
    \item Ψ-dimensions must be exhaustive for behavioral prediction
    \item Modulation functions must be empirically calibratable
\end{itemize}

\textbf{Lieferobjekte (Deliverables):}
\begin{itemize}[nosep]
    \item Complete Ψ-Taxonomy (Table~\ref{tab:when-psi})
    \item Axiom System PSI-1 through PSI-10
    \item Context Modulation Formula (Equation~\ref{eq:when-modulation})
    \item MACRO-MESO-MICRO specification
\end{itemize}

\end{tcolorbox}

% =============================================================================
% CHAPTER
% =============================================================================

\chapter{CORE-WHEN: Context Framework (Ψ-Dimensions)}
\label{app:core-when}

\begin{abstract}
This appendix answers the fourth fundamental question of the Evidence-Based Framework:
\textit{``When does context matter?''} We establish that all behavioral parameters
are context-dependent, formalized through 8 Ψ-dimensions: Institutional ($\Psi_I$),
Social ($\Psi_S$), Cognitive ($\Psi_C$), Cultural ($\Psi_K$), Economic ($\Psi_E$),
Temporal ($\Psi_T$), Material ($\Psi_M$), and Physical ($\Psi_F$). The axiom system
PSI-1 through PSI-10 specifies how context modulates utility weights, complementarity
values, and behavioral parameters. The MACRO-MESO-MICRO hierarchy provides a
structured approach to context specification.
\end{abstract}

% -----------------------------------------------------------------------------
% QUICK REFERENCE BOX
% -----------------------------------------------------------------------------

\begin{tcolorbox}[colback=blue!5!white, colframe=blue!75!black,
    title=Quick Reference: CORE-WHEN]

\textbf{Core Question:} When does context matter? How do situations modulate behavior?

\textbf{Key Formula:}
\begin{equation}
\theta_{\text{eff}} = \theta_{\text{base}} \times \prod_{k \in \{I,S,C,K,E,T,M,F\}} m_k(\Psi_k)
\label{eq:when-modulation}
\end{equation}

\textbf{The 8 Ψ-Dimensions:}
\begin{itemize}[nosep]
\item $\Psi_I$: \textbf{I}nstitutional --- Rules, regulations, defaults
\item $\Psi_S$: \textbf{S}ocial --- Who is present, social norms
\item $\Psi_C$: \textbf{C}ognitive --- Mental state, attention, load
\item $\Psi_K$: \textbf{K}ultural --- Values, beliefs, traditions
\item $\Psi_E$: \textbf{E}conomic --- Resources, scarcity, incentives
\item $\Psi_T$: \textbf{T}emporal --- Time, urgency, life phase
\item $\Psi_M$: \textbf{M}aterial --- Technology, tools, infrastructure
\item $\Psi_F$: Physical (\textbf{F}isico) --- Location, environment
\end{itemize}

\textbf{Cross-References:}
AAA (CORE-WHO), C (CORE-WHAT), B (CORE-HOW), AU (CORE-AWARE), CAL (METHOD-LLMMC), CP (DOMAIN-CPS: Intransparenz \& Dynamik)

\end{tcolorbox}


%%%%%%%%%%%%%%%%%%%%%%%%%%%%%%%%%%%%%%%%%%%%%%%%%%%%%%%%%%%%%%%%%%%%%%%%%%%%%%%
\section{The Fundamental Question}
\label{sec:when-fundamental}
%%%%%%%%%%%%%%%%%%%%%%%%%%%%%%%%%%%%%%%%%%%%%%%%%%%%%%%%%%%%%%%%%%%%%%%%%%%%%%%

\subsection{Why "When" Matters}

Standard economics assumes stable preferences and context-independent parameters.
But the same person makes different choices in different contexts:

\begin{itemize}
\item Loss aversion varies with stakes and framing ($\Psi_I$, $\Psi_E$)
\item Social preferences depend on who is observing ($\Psi_S$)
\item Time discounting varies with cognitive load ($\Psi_C$)
\item Risk preferences differ across cultures ($\Psi_K$)
\item Defaults matter more under time pressure ($\Psi_T$)
\end{itemize}

\begin{tcolorbox}[colback=yellow!5!white, colframe=yellow!75!black,
    title=Context as Multiplier]
\textbf{The Central Thesis:}
\[
\text{Same Question} + \text{Different Context} = \text{Different Answer}
\]

Context is not noise to be controlled away --- it is the \textbf{primary determinant}
of whether an intervention works or fails.
\end{tcolorbox}

\subsection{The Core Problem}

Ignoring context leads to:

\begin{enumerate}
\item \textbf{Failed Replications}: Interventions that work in one context fail in another
\item \textbf{WEIRD Bias}: Lab findings from Western samples don't generalize
\item \textbf{Prediction Failure}: Models without context systematically under-perform
\end{enumerate}

\begin{tcolorbox}[colback=red!5!white, colframe=red!75!black,
    title=The Fundamental Question]
\textbf{When does context matter? What are the dimensions of situational variation that modulate behavioral parameters?}
\end{tcolorbox}


%%%%%%%%%%%%%%%%%%%%%%%%%%%%%%%%%%%%%%%%%%%%%%%%%%%%%%%%%%%%%%%%%%%%%%%%%%%%%%%
\section{Core Theory: The 8 Ψ-Dimensions}
\label{sec:when-dimensions}
%%%%%%%%%%%%%%%%%%%%%%%%%%%%%%%%%%%%%%%%%%%%%%%%%%%%%%%%%%%%%%%%%%%%%%%%%%%%%%%

\subsection{Overview}

\begin{table}[htbp]
\centering
\caption{The 8 Ψ-Dimensions of Context}
\label{tab:when-psi}
\renewcommand{\arraystretch}{1.4}
\begin{tabular}{@{}clp{4.5cm}l@{}}
\toprule
\textbf{Code} & \textbf{Dimension} & \textbf{Definition} & \textbf{Example Variables} \\
\midrule
$\Psi_I$ & Institutional & Rules, regulations, defaults, choice architecture & Opt-in vs. opt-out, legal constraints \\
$\Psi_S$ & Social & Social presence, norms, observation, peers & Alone vs. observed, group size \\
$\Psi_C$ & Cognitive & Mental state, attention, load, motivation & Fatigue, stress, distraction \\
$\Psi_K$ & Cultural & Values, beliefs, traditions, worldviews & Individualism, trust, religion \\
$\Psi_E$ & Economic & Resources, scarcity, stakes, incentives & Income, wealth, budget constraints \\
$\Psi_T$ & Temporal & Time, urgency, deadline, life phase & Morning vs. evening, age, career stage \\
$\Psi_M$ & Material & Technology, tools, interface, infrastructure & Digital vs. paper, mobile vs. desktop \\
$\Psi_F$ & Physical & Location, environment, weather, space & Home vs. office, urban vs. rural \\
\bottomrule
\end{tabular}
\end{table}

\subsection{Detailed Dimension Specifications}

\subsubsection{$\Psi_I$: Institutional Context}

\begin{definition}[Institutional Context]
\textbf{Institutional context} $\Psi_I$ captures the rule structure governing choices:
regulations, defaults, framing, and choice architecture.
\end{definition}

\textbf{Key Variables:}
\begin{itemize}[nosep]
\item Default setting (opt-in vs. opt-out)
\item Legal constraints and permissions
\item Framing (gain vs. loss)
\item Information disclosure requirements
\item Enforcement probability
\end{itemize}

\textbf{Behavioral Implication:}
Default compliance rate is approximately 90\% under opt-out vs. 10\% under opt-in
(\citet{johnson2003defaults}).

\subsubsection{$\Psi_S$: Social Context}

\begin{definition}[Social Context]
\textbf{Social context} $\Psi_S$ captures the social environment:
who is present, social norms, observation, and reference groups.
\end{definition}

\textbf{Key Variables:}
\begin{itemize}[nosep]
\item Presence of others (alone vs. observed)
\item Relationship type (strangers, friends, family, authority)
\item Group size and composition
\item Descriptive norms (what others do)
\item Injunctive norms (what others approve)
\item \textbf{Network structure} $\rho \in [0,1]$ (mixing vs. segregation)
\end{itemize}

\textbf{Network Structure Parameter $\rho$:}
\begin{itemize}[nosep]
\item $\rho = 0$: Complete mixing (heterogeneous interactions)
\item $\rho = 1$: Complete segregation (homophilous interactions)
\item Higher $\rho$ = narratives spread within homogeneous groups
\item Lower $\rho$ = exposure to diverse perspectives, higher prosocial behavior
\end{itemize}

\textbf{Behavioral Implication:}
Social observation increases prosocial behavior by approximately 25\%
(\citet{ariely2009doing}). Network mixing (lower $\rho$) increases moral behavior
by exposing agents to prosocial types (\citet{benabou_2018_narratives}).

\subsubsection{$\Psi_C$: Cognitive Context}

\begin{definition}[Cognitive Context]
\textbf{Cognitive context} $\Psi_C$ captures the mental state of the decision-maker:
attention, cognitive load, fatigue, and arousal.
\end{definition}

\textbf{Key Variables:}
\begin{itemize}[nosep]
\item Cognitive load (low/medium/high)
\item Fatigue level
\item Stress and arousal
\item Attention availability
\item Decision fatigue
\end{itemize}

\textbf{Behavioral Implication:}
High cognitive load increases reliance on heuristics and defaults
(\citet{kahneman2011thinking}).

\subsubsection{$\Psi_K$: Cultural Context}

\begin{definition}[Cultural Context]
\textbf{Cultural context} $\Psi_K$ captures the cultural environment:
values, beliefs, traditions, and worldviews.
\end{definition}

\textbf{Key Variables:}
\begin{itemize}[nosep]
\item Individualism vs. collectivism
\item Power distance
\item Uncertainty avoidance
\item Institutional trust
\item Religious orientation
\end{itemize}

\textbf{Behavioral Implication:}
Loss aversion is approximately 40\% higher in collectivist cultures
(\citet{wang2017cultural}).

\subsubsection{$\Psi_E$: Economic Context}

\begin{definition}[Economic Context]
\textbf{Economic context} $\Psi_E$ captures resource availability:
income, wealth, scarcity, and economic incentives.
\end{definition}

\textbf{Key Variables:}
\begin{itemize}[nosep]
\item Income level (absolute and relative)
\item Wealth and assets
\item Scarcity (bandwidth tax)
\item Stakes (high vs. low)
\item Price sensitivity
\end{itemize}

\textbf{Behavioral Implication:}
Scarcity narrows bandwidth, reducing cognitive capacity for non-urgent decisions
(\citet{mullainathan2013scarcity}).

\subsubsection{$\Psi_T$: Temporal Context}

\begin{definition}[Temporal Context]
\textbf{Temporal context} $\Psi_T$ captures time-related factors:
urgency, deadlines, time of day, life phase, and \textit{time as scarce resource}.
\end{definition}

\textbf{Theoretical Foundation: Time Allocation Theory}

The foundational framework for understanding temporal context comes from
\citet{PAP-becker1965theory}, who established time as an economic resource with
opportunity cost. The key insight is that time is finite and must be allocated
across activities:
\begin{equation}
T = T_{\text{market}} + T_{\text{household}} + T_{\text{leisure}}
\end{equation}

The \textbf{shadow price of time} equals the marginal wage rate $w$, implying
that the full price of any activity includes both monetary cost and time cost:
\begin{equation}
\pi_i = p_i + w \cdot t_i
\end{equation}

This framework (formalized as MS-TA-001 in the Theory Catalog, CAT-17) explains
phenomena from female labor supply growth to digital platform competition
(see Appendix PT, CAS-TA-001, CAS-853).

\textbf{Key Variables:}
\begin{itemize}[nosep]
\item Time pressure (deadline proximity)
\item Time of day (circadian effects)
\item Day of week
\item Life phase (age, career stage)
\item Planning horizon
\item Time budget constraint (Becker)
\item Opportunity cost of time ($w$)
\end{itemize}

\textbf{Behavioral Implication:}
Willpower and self-control deplete throughout the day
(\citet{baumeister2007strength}). Additionally, rising wages increase
the opportunity cost of time, leading to substitution from time-intensive
to goods-intensive activities (\citet{PAP-becker1965theory}).

\subsubsection{$\Psi_M$: Material/Technological Context}

\begin{definition}[Material Context]
\textbf{Material context} $\Psi_M$ captures the technological and physical infrastructure:
tools, interfaces, and material affordances.
\end{definition}

\textbf{Key Variables:}
\begin{itemize}[nosep]
\item Digital vs. physical interface
\item Mobile vs. desktop
\item Payment method (cash vs. card vs. digital)
\item Tool availability
\item Interface design
\end{itemize}

\textbf{Behavioral Implication:}
Mobile interfaces reduce attention span and increase impulsive choices
(\citet{ghose2013mobile}).

\subsubsection{$\Psi_F$: Physical Context}

\begin{definition}[Physical Context]
\textbf{Physical context} $\Psi_F$ captures the physical environment:
location, weather, temperature, and spatial configuration.
\end{definition}

\textbf{Key Variables:}
\begin{itemize}[nosep]
\item Location (home, office, public space)
\item Weather and temperature
\item Noise level
\item Spatial design (open vs. private)
\item Urban vs. rural
\end{itemize}

\textbf{Behavioral Implication:}
Warm temperatures increase risk-taking and social affiliation
(\citet{huang2014warm}).


%%%%%%%%%%%%%%%%%%%%%%%%%%%%%%%%%%%%%%%%%%%%%%%%%%%%%%%%%%%%%%%%%%%%%%%%%%%%%%%
\section{Axioms and Formal Definitions}
\label{sec:when-axioms}
%%%%%%%%%%%%%%%%%%%%%%%%%%%%%%%%%%%%%%%%%%%%%%%%%%%%%%%%%%%%%%%%%%%%%%%%%%%%%%%

% PSI-1
\begin{tcolorbox}[colback=gray!5!white, colframe=gray!75!black,
    title=Axiom PSI-1: Universal Context Dependence]
\textbf{Statement:}
\begin{equation}
\forall \theta \in \Theta_{\text{behavioral}}: \theta = \theta(\Psi)
\end{equation}

\textbf{Interpretation:} All behavioral parameters are context-dependent.
There are no truly context-free parameters.

\textbf{Implication:} Parameters reported without context specification are incomplete.
\end{tcolorbox}

% PSI-2
\begin{tcolorbox}[colback=gray!5!white, colframe=gray!75!black,
    title=Axiom PSI-2: Multiplicative Modulation]
\textbf{Statement:}
\begin{equation}
\theta_{\text{eff}} = \theta_{\text{base}} \times \prod_{k} m_k(\Psi_k)
\label{eq:when-modulation-axiom}
\end{equation}

\textbf{Interpretation:} Context modulation is multiplicative (not additive).
Each Ψ-dimension applies a multiplier to the base parameter.

\textbf{Implication:} Extreme contexts can dramatically amplify or suppress effects.
\end{tcolorbox}

% PSI-3
\begin{tcolorbox}[colback=gray!5!white, colframe=gray!75!black,
    title=Axiom PSI-3: Dimension Completeness]
\textbf{Statement:}
\begin{equation}
\Psi = (\Psi_I, \Psi_S, \Psi_C, \Psi_K, \Psi_E, \Psi_T, \Psi_M, \Psi_F) \text{ is exhaustive}
\end{equation}

\textbf{Interpretation:} The 8 Ψ-dimensions capture all behaviorally relevant context.
Any context factor maps to at least one dimension.

\textbf{Implication:} New context factors are classified within existing dimensions.
\end{tcolorbox}

% PSI-4
\begin{tcolorbox}[colback=gray!5!white, colframe=gray!75!black,
    title=Axiom PSI-4: Dimension Interactions]
\textbf{Statement:}
\begin{equation}
\exists \; k, k': \frac{\partial^2 m}{\partial \Psi_k \partial \Psi_{k'}} \neq 0
\end{equation}

\textbf{Interpretation:} Ψ-dimensions can interact. The effect of institutional context
may depend on cultural context (e.g., defaults work differently in high-trust cultures).

\textbf{Cross-Reference:} Interaction structure parallels CORE-HOW (B) complementarity.
\end{tcolorbox}

% PSI-5
\begin{tcolorbox}[colback=gray!5!white, colframe=gray!75!black,
    title=Axiom PSI-5: Hierarchy (MACRO-MESO-MICRO)]
\textbf{Statement:}
\begin{equation}
\Psi = \Psi_{\text{MACRO}}(\text{country}) \times \Psi_{\text{MESO}}(\text{sector}) \times \Psi_{\text{MICRO}}(\text{situation})
\end{equation}

\textbf{Interpretation:} Context operates at three hierarchical levels:
MACRO (country/market), MESO (industry/organization), MICRO (immediate situation).

\textbf{Implication:} Context analysis proceeds top-down: country → sector → situation.
\end{tcolorbox}

% PSI-6 through PSI-10 (abbreviated)
\begin{tcolorbox}[colback=gray!5!white, colframe=gray!75!black,
    title=Axioms PSI-6 to PSI-10: Modulation Properties]

\textbf{PSI-6 (Bounded Modulation):} Multipliers are bounded: $m_k \in [m_{\min}, m_{\max}]$

\textbf{PSI-7 (Asymmetric Effects):} Modulation can be asymmetric (gain vs. loss contexts)

\textbf{PSI-8 (Temporal Dynamics):} Context effects have decay and reinforcement dynamics

\textbf{PSI-9 (Individual Sensitivity):} Individuals vary in sensitivity to context

\textbf{PSI-10 (Measurement Requirement):} Each Ψ-dimension must be operationalizable
\end{tcolorbox}


%%%%%%%%%%%%%%%%%%%%%%%%%%%%%%%%%%%%%%%%%%%%%%%%%%%%%%%%%%%%%%%%%%%%%%%%%%%%%%%
\section{The MACRO-MESO-MICRO Hierarchy}
\label{sec:when-hierarchy}
%%%%%%%%%%%%%%%%%%%%%%%%%%%%%%%%%%%%%%%%%%%%%%%%%%%%%%%%%%%%%%%%%%%%%%%%%%%%%%%

\subsection{Three-Level Structure}

\begin{table}[htbp]
\centering
\caption{MACRO-MESO-MICRO Context Hierarchy}
\label{tab:when-hierarchy}
\renewcommand{\arraystretch}{1.3}
\begin{tabular}{@{}clll@{}}
\toprule
\textbf{Level} & \textbf{Name} & \textbf{Scope} & \textbf{Data Source} \\
\midrule
1 & MACRO & Country/Market & BCM2\_04\_KON\_*.yaml \\
2 & MESO & Industry/Organization & clients/*, customers/* \\
3 & MICRO & Immediate Situation & context-dimensions.yaml \\
\bottomrule
\end{tabular}
\end{table}

\subsection{Cascade Calculation}

The final context parameter is:

\begin{equation}
\Psi_{\text{final}} = \Psi_{\text{MACRO}} \times \Psi_{\text{MESO}} \times \Psi_{\text{MICRO}}
\end{equation}

\textbf{Example:}
\begin{align}
\lambda_{\text{final}} &= \lambda_{\text{base}} \times m_{\text{CH}} \times m_{\text{BFE}} \times m_{\text{situation}} \\
&= 2.0 \times 1.05 \times 1.2 \times 1.3 \\
&= 3.28
\end{align}


%%%%%%%%%%%%%%%%%%%%%%%%%%%%%%%%%%%%%%%%%%%%%%%%%%%%%%%%%%%%%%%%%%%%%%%%%%%%%%%
\section{Context Modulation Rules}
\label{sec:when-modulation}
%%%%%%%%%%%%%%%%%%%%%%%%%%%%%%%%%%%%%%%%%%%%%%%%%%%%%%%%%%%%%%%%%%%%%%%%%%%%%%%

\subsection{Default Modulation Functions}

For each Ψ-dimension, the default modulation function is:

\begin{equation}
m_k(\Psi_k) = 1 + \beta_k \cdot (\Psi_k - \Psi_k^{\text{neutral}})
\end{equation}

where:
\begin{itemize}
\item $\beta_k$ = sensitivity coefficient for dimension $k$
\item $\Psi_k^{\text{neutral}}$ = neutral reference point
\end{itemize}

\subsection{Known Modulation Effects}

\begin{table}[htbp]
\centering
\caption{Empirical Modulation Effects}
\label{tab:when-effects}
\renewcommand{\arraystretch}{1.3}
\begin{tabular}{@{}llcc@{}}
\toprule
\textbf{Dimension} & \textbf{Variable} & \textbf{Effect on λ} & \textbf{Source} \\
\midrule
$\Psi_I$ & Opt-out default & $\times 1.5$ & Johnson \& Goldstein \\
$\Psi_S$ & Social observation & $\times 1.25$ & Ariely \\
$\Psi_S$ & Network mixing (low $\rho$) & $\times 1.15$ & Bénabou et al. \\
$\Psi_C$ & High cognitive load & $\times 1.3$ & Kahneman \\
$\Psi_K$ & Collectivism (high) & $\times 1.4$ & Wang et al. \\
$\Psi_E$ & High stakes & $\times 1.2$ & Camerer \\
$\Psi_T$ & Deadline pressure & $\times 1.3$ & Ariely \& Wertenbroch \\
\bottomrule
\end{tabular}
\end{table}


%%%%%%%%%%%%%%%%%%%%%%%%%%%%%%%%%%%%%%%%%%%%%%%%%%%%%%%%%%%%%%%%%%%%%%%%%%%%%%%
\section{Integration with Other CORE Appendices}
\label{sec:when-integration}
%%%%%%%%%%%%%%%%%%%%%%%%%%%%%%%%%%%%%%%%%%%%%%%%%%%%%%%%%%%%%%%%%%%%%%%%%%%%%%%

\begin{tcolorbox}[colback=yellow!5!white, colframe=yellow!75!black,
    title=Integration: CORE-WHEN $\leftrightarrow$ CORE-WHO]
\textbf{What WHEN provides to WHO:}
\begin{itemize}[nosep]
\item Context-dependent welfare weights: $\omega_L(\Psi)$
\item Level-specific context relevance
\end{itemize}

\textbf{What WHEN requires from WHO:}
\begin{itemize}[nosep]
\item Level structure for context application
\item Entity specification per level
\end{itemize}
\end{tcolorbox}

\begin{tcolorbox}[colback=yellow!5!white, colframe=yellow!75!black,
    title=Integration: CORE-WHEN $\leftrightarrow$ CORE-WHAT]
\textbf{What WHEN provides to WHAT:}
\begin{itemize}[nosep]
\item Context-dependent dimension weights: $\alpha_d(\Psi)$
\item When dimension importance shifts
\end{itemize}

\textbf{What WHEN requires from WHAT:}
\begin{itemize}[nosep]
\item Dimension structure (FEPSDE)
\item Baseline dimension weights
\end{itemize}
\end{tcolorbox}

\begin{tcolorbox}[colback=yellow!5!white, colframe=yellow!75!black,
    title=Integration: CORE-WHEN $\leftrightarrow$ CORE-HOW]
\textbf{What WHEN provides to HOW:}
\begin{itemize}[nosep]
\item Context-dependent complementarity: $\gamma_{ij}(\Psi)$
\item When crowding-out vs. synergy shifts
\end{itemize}

\textbf{What WHEN requires from HOW:}
\begin{itemize}[nosep]
\item Baseline γ-matrix
\item Interaction structure
\end{itemize}
\end{tcolorbox}

\begin{tcolorbox}[colback=green!5!white, colframe=green!75!black,
    title=Application Domain: Complex Problem Solving (CP)]
\textbf{CPS leverages three WHEN-specific parameters:}
\begin{itemize}[nosep]
\item \textbf{Intransparenz ($\eta$):} Initial and ongoing opacity of the problem system
\item \textbf{Dynamik ($dX/dt$):} Autonomous change rate of problem variables
\item \textbf{Feedback-Delay ($\delta$):} Time lag between action and observable consequence
\end{itemize}

\textbf{Main Finding (Dörner 1983, 1989):} Intransparenz is the dominant lever ($\sim 42\%$ influence on CPS success).

\textbf{Cross-Reference:} See Appendix CP (DOMAIN-CPS) for the complete 10C-CPS mapping and intervention recommendations.
\end{tcolorbox}


%%%%%%%%%%%%%%%%%%%%%%%%%%%%%%%%%%%%%%%%%%%%%%%%%%%%%%%%%%%%%%%%%%%%%%%%%%%%%%%
\section{Worked Example}
\label{sec:when-example}
%%%%%%%%%%%%%%%%%%%%%%%%%%%%%%%%%%%%%%%%%%%%%%%%%%%%%%%%%%%%%%%%%%%%%%%%%%%%%%%

\subsection{Example: Retirement Savings Intervention}

\paragraph{Setup.}
Design an intervention to increase retirement savings in Switzerland (MACRO)
for BFE employees (MESO) during annual benefits enrollment (MICRO).

\paragraph{Context Analysis.}

\textbf{MACRO (Switzerland):}
\begin{itemize}[nosep]
\item $\Psi_K$: High institutional trust (m = 1.1)
\item $\Psi_I$: Strong regulatory framework (m = 1.0)
\item Base $\lambda_{\text{CH}} = 2.1$
\end{itemize}

\textbf{MESO (BFE Energy):}
\begin{itemize}[nosep]
\item $\Psi_E$: Above-average income (m = 0.9 for loss aversion)
\item $\Psi_S$: Professional culture (m = 1.0)
\end{itemize}

\textbf{MICRO (Enrollment Period):}
\begin{itemize}[nosep]
\item $\Psi_T$: Deadline pressure (m = 1.3)
\item $\Psi_C$: Medium cognitive load (m = 1.1)
\item $\Psi_I$: Opt-out default available (m = 1.5)
\end{itemize}

\paragraph{Effective Parameters.}
\begin{align}
\lambda_{\text{eff}} &= 2.1 \times 1.1 \times 0.9 \times 1.3 \times 1.1 = 3.0 \\
\text{Default effect} &= 1.5 \rightarrow \text{Opt-out default recommended}
\end{align}

\paragraph{Recommendation.}
Use opt-out enrollment with loss-framed communication
("Don't miss out on CHF X in employer matching").


%%%%%%%%%%%%%%%%%%%%%%%%%%%%%%%%%%%%%%%%%%%%%%%%%%%%%%%%%%%%%%%%%%%%%%%%%%%%%%%
\section{Implications for Practice}
\label{sec:when-practice}
%%%%%%%%%%%%%%%%%%%%%%%%%%%%%%%%%%%%%%%%%%%%%%%%%%%%%%%%%%%%%%%%%%%%%%%%%%%%%%%

\begin{tcolorbox}[colback=green!5!white, colframe=green!75!black,
    title=Practical Implications]
\begin{enumerate}
\item \textbf{Always Specify Context:} No parameter or intervention recommendation
      is valid without context specification
\item \textbf{Use MACRO-MESO-MICRO:} Start with country, then sector, then situation
\item \textbf{Context Before Intervention:} Analyze context before designing interventions
\item \textbf{Exploit Favorable Contexts:} Time interventions for favorable
      Ψ-configurations (e.g., enrollment periods, life transitions)
\end{enumerate}
\end{tcolorbox}


%%%%%%%%%%%%%%%%%%%%%%%%%%%%%%%%%%%%%%%%%%%%%%%%%%%%%%%%%%%%%%%%%%%%%%%%%%%%%%%
\section{Summary}
\label{sec:when-summary}
%%%%%%%%%%%%%%%%%%%%%%%%%%%%%%%%%%%%%%%%%%%%%%%%%%%%%%%%%%%%%%%%%%%%%%%%%%%%%%%

\begin{tcolorbox}[colback=gray!5!white, colframe=gray!75!black,
    title=Appendix V Summary: CORE-WHEN]

\textbf{Core Contribution:}
Establishes that all behavioral parameters are context-dependent through
8 Ψ-dimensions, with multiplicative modulation.

\textbf{Key Formula:}
\begin{equation}
\theta_{\text{eff}} = \theta_{\text{base}} \times \prod_{k \in \{I,S,C,K,E,T,M,F\}} m_k(\Psi_k)
\end{equation}

\textbf{Key Insights:}
\begin{itemize}[nosep]
\item Context is not noise; it is the primary determinant of intervention success
\item 8 Ψ-dimensions are exhaustive for behavioral prediction
\item MACRO-MESO-MICRO provides structured context analysis
\item Same question + different context = different answer
\end{itemize}

\textbf{Next Steps:}
See Appendix BBB (CORE-WHERE) for parameter estimation methods.
See Appendix CAL (METHOD-LLMMC) for calibration procedures.

\end{tcolorbox}


%%%%%%%%%%%%%%%%%%%%%%%%%%%%%%%%%%%%%%%%%%%%%%%%%%%%%%%%%%%%%%%%%%%%%%%%%%%%%%%
\section{Glossary of Symbols}
\label{sec:when-glossary}
%%%%%%%%%%%%%%%%%%%%%%%%%%%%%%%%%%%%%%%%%%%%%%%%%%%%%%%%%%%%%%%%%%%%%%%%%%%%%%%

\begin{table}[htbp]
\centering
\caption{Symbols Introduced in Appendix V}
\label{tab:when-symbols}
\renewcommand{\arraystretch}{1.3}
\begin{tabular}{@{}clp{5cm}c@{}}
\toprule
\textbf{Symbol} & \textbf{Name} & \textbf{Definition} & \textbf{Status} \\
\midrule
$\Psi$ & Context Vector & $(\Psi_I, \Psi_S, \Psi_C, \Psi_K, \Psi_E, \Psi_T, \Psi_M, \Psi_F)$ & NEW \\
$\Psi_k$ & Context Dimension & Single dimension of context & NEW \\
$m_k(\cdot)$ & Modulation Function & Maps context to multiplier & NEW \\
$\theta_{\text{eff}}$ & Effective Parameter & Context-adjusted parameter & NEW \\
$\theta_{\text{base}}$ & Base Parameter & Context-free reference value & NEW \\
$\rho$ & Network Persistence & Social network segregation $\in [0,1]$ & NEW \\
\bottomrule
\end{tabular}
\end{table}


%%%%%%%%%%%%%%%%%%%%%%%%%%%%%%%%%%%%%%%%%%%%%%%%%%%%%%%%%%%%%%%%%%%%%%%%%%%%%%%
\section{Critical Foundations and Objections}
\label{sec:when-critical}
%%%%%%%%%%%%%%%%%%%%%%%%%%%%%%%%%%%%%%%%%%%%%%%%%%%%%%%%%%%%%%%%%%%%%%%%%%%%%%%

\subsection{Objection 1: Infinite Regress}

\begin{tcolorbox}[colback=red!5!white, colframe=red!50!black,
    title=Objection: Context Depends on Context]
\textit{``If all parameters are context-dependent, then modulation coefficients
are also context-dependent. This leads to infinite regress.''}
\end{tcolorbox}

\textbf{Response:}
The framework is hierarchical, not infinitely recursive. MACRO context is treated
as fixed within an analysis; MESO and MICRO modulate within MACRO. Modulation
coefficients $m_k$ are empirically calibrated constants at each level.

\subsection{Objection 2: Overfitting Risk}

\begin{tcolorbox}[colback=red!5!white, colframe=red!50!black,
    title=Objection: Too Many Free Parameters]
\textit{``With 8 dimensions and multiple variables per dimension, any outcome
can be explained post-hoc. The framework is unfalsifiable.''}
\end{tcolorbox}

\textbf{Response:}
PSI-10 (Measurement Requirement) constrains dimensions to operationalizable factors.
Empirical calibration (Appendix BBB) provides default values from meta-analyses.
Out-of-sample prediction is the test, not in-sample fit.


%%%%%%%%%%%%%%%%%%%%%%%%%%%%%%%%%%%%%%%%%%%%%%%%%%%%%%%%%%%%%%%%%%%%%%%%%%%%%%%
\section{Open Issues and Research Agenda}
\label{sec:when-open}
%%%%%%%%%%%%%%%%%%%%%%%%%%%%%%%%%%%%%%%%%%%%%%%%%%%%%%%%%%%%%%%%%%%%%%%%%%%%%%%

\begin{table}[htbp]
\centering
\caption{Research Roadmap for CORE-WHEN}
\label{tab:when-roadmap}
\renewcommand{\arraystretch}{1.3}
\begin{tabular}{@{}clccl@{}}
\toprule
\textbf{\#} & \textbf{Issue} & \textbf{Priority} & \textbf{Status} & \textbf{Requirements} \\
\midrule
1 & Complete $m_k$ calibration & HIGH & Partial & Meta-analysis \\
2 & Ψ-dimension interactions & HIGH & Open & Factorial designs \\
3 & MICRO context measurement & MEDIUM & Ongoing & ESM methods \\
4 & Temporal dynamics of Ψ & MEDIUM & Open & Panel data \\
5 & AI-assisted context inference & LOW & Exploratory & ML methods \\
\bottomrule
\end{tabular}
\end{table}


%%%%%%%%%%%%%%%%%%%%%%%%%%%%%%%%%%%%%%%%%%%%%%%%%%%%%%%%%%%%%%%%%%%%%%%%%%%%%%%
\section{References}
\label{sec:when-references}
%%%%%%%%%%%%%%%%%%%%%%%%%%%%%%%%%%%%%%%%%%%%%%%%%%%%%%%%%%%%%%%%%%%%%%%%%%%%%%%

\begin{tcolorbox}[colback=blue!5!white, colframe=blue!75!black]
\textbf{Master Bibliography:} See \texttt{bcm\_master.bib}
\end{tcolorbox}

\subsection{Key References}

\begin{itemize}[nosep]
\item \textbf{Defaults:} \citet{johnson2003defaults}
\item \textbf{Social Context:} \citet{ariely2009doing}
\item \textbf{Network Structure:} \citet{benabou_2018_narratives} --- Network parameter $\rho$ and moral transmission
\item \textbf{Cognitive Load:} \citet{kahneman2011thinking}
\item \textbf{Cultural Context:} \citet{wang2017cultural}, \citet{hofstede2010cultures}
\item \textbf{Scarcity:} \citet{mullainathan2013scarcity}
\end{itemize}

\nocite{bcm_master}

%%%%%%%%%%%%%%%%%%%%%%%%%%%%%%%%%%%%%%%%%%%%%%%%%%%%%%%%%%%%%%%%%%%%%%%%%%%%%%%
% FOOTER
%%%%%%%%%%%%%%%%%%%%%%%%%%%%%%%%%%%%%%%%%%%%%%%%%%%%%%%%%%%%%%%%%%%%%%%%%%%%%%%

\vspace{1em}
\hrule
\vspace{0.5em}

\noindent\textit{Cross-references:}
Appendix GLS (Glossary),
Appendix AAA (CORE-WHO),
Appendix C (CORE-WHAT),
Appendix B (CORE-HOW),
Appendix BBB (CORE-WHERE),
Appendix AU (CORE-AWARE: Moral Narratives \& Network Transmission),
Appendix RB (LIT-BENABOU: Network Parameter $\rho$),
Appendix CAL (METHOD-LLMMC).

\noindent\textit{Master Bibliography:} \texttt{bcm\_master.bib}

\noindent\textit{Note:} This appendix replaces former CTW references.

% =============================================================================
% END OF APPENDIX V
% =============================================================================
